% =============================================================================
\section{Introduction to \bert{}}
\label{sec:intro_bert}
% =============================================================================

\subsection{From Transformer to \bert{}}

The Transformer's encoder already attends bidirectionally: in each
self-attention head, every token can query every other token regardless of
direction. However, the original Transformer was designed as an encoder--decoder
for sequence-to-sequence tasks. \textbf{\bert{}}~\cite{devlin2019bert} (Devlin
et al., 2018) makes a deliberate choice: \emph{keep only the encoder stack} and
pre-train it for language understanding rather than generation.

\begin{figure}[H]
  \centering
  \includegraphics[width=\textwidth]{comparison.png}
  \caption{Architectural comparison. The original Transformer uses both
    encoder and decoder for translation. \bert{} retains only the encoder stack,
    using the \cls{} token representation as input to task-specific classifier
    heads.}
  \label{fig:bert_vs_transformer}
\end{figure}

Table~\ref{tab:transformer_vs_bert} contrasts the two architectures.

\begin{table}[H]
\centering
\caption{Original Transformer vs.\ \bert{}.}
\label{tab:transformer_vs_bert}
\begin{tabularx}{0.92\textwidth}{lXX}
  \toprule
  \textbf{Aspect}
    & \textbf{Original Transformer}~\cite{vaswani2017attention}
    & \textbf{\bert{}}~\cite{devlin2019bert} \\
  \midrule
  Architecture
    & Encoder + Decoder
    & Encoder only ($N$ stacked layers) \\
  Primary task
    & Sequence-to-sequence (translation)
    & Representation learning / understanding \\
  Context
    & Decoder is left-to-right (masked)
    & Fully bidirectional at every layer \\
  Pre-training
    & None (direct supervision)
    & Masked LM + Next Sentence Prediction \\
  Usage
    & Train end-to-end per task
    & Pre-train once, fine-tune per task \\
  \bottomrule
\end{tabularx}
\end{table}

\begin{remark}
  \textit{``The Transformer learns to \textbf{translate}. \bert{} learns to
  \textbf{understand}.''}
\end{remark}

\subsection{What Makes \bert{} Truly Bidirectional?}

This is a subtle but crucial point. ELMo~\cite{peters2018elmo} was also called
``bidirectional,'' but its forward and backward LSTMs were trained with separate
language model objectives and then concatenated. \bert{}, by contrast, uses a
\emph{single unified} Transformer encoder in which every attention head
simultaneously sees all tokens from both directions, and every layer integrates
left and right context jointly. This \emph{deep} bidirectionality --- as opposed
to shallow concatenation --- is what makes \bert{}'s representations
qualitatively different.